\chapter{Introduction}
\label{ch:introduction}


\section{Context and Motivation}

Sepsis is a life-threatening condition caused by the body's extreme response to
infection \cite{singer2016sepsis3}. It kills about eleven million people
every year, nearly one in five deaths worldwide
\cite{rudd2020global}. In the ICU, every hour of delayed treatment raises
the risk of death \cite{seymour2017time}. An early warning system that
detects sepsis before doctors recognise it could save lives by starting
treatment sooner.

This potential has driven major research investment. Wang et al.\
\cite{wang2025srpm} reviewed ninety-one published real-time sepsis
prediction models and found a median internal AUROC of 0.811, suggesting
strong performance. But the same review found that the clinical Utility Score
collapsed from $+0.381$ internally to $-0.164$ externally. A negative score
means the model's alerts cause more harm than benefit. AUROC did not detect
this collapse because it only measures whether a model can rank patients
correctly, not whether its alerts are timely or useful. This gap between
apparent accuracy and real-world usefulness has also been observed in
deployment: Wong et al.\ \cite{wong2026epic} validated the Epic Sepsis
Model v2 prospectively across 227,091 encounters at four US health systems and
found encounter-level AUROC between 0.82 and 0.92 alongside positive predictive
value of 0.13 to 0.26: 21 to 35 alerts investigated for every sepsis case
identified.

The key question is therefore not ``which model is best?'' but something more
fundamental: \textit{how much of a model's reported performance comes from the
model itself, and how much comes from how the study was designed: how sepsis
was defined, how data was split, and which metric was reported?}

Cohen et al.\ \cite{cohen2024subtle} showed on MIMIC-III that varying only
the Sepsis-3 definition changed the identified cohort from 867 to 2,178
patients, and that the label choice affected performance as much as the model
choice. Kamran et al.\ \cite{kamran2024evaluation} showed that published
models partly learn to predict when clinicians start treatment rather than when
patients deteriorate. Together, these findings suggest that much of the reported
performance may come from the researcher's design choices, not from the AI model.

This thesis tests that hypothesis directly. Instead of fixing the label, the
split, and the metric and asking which model wins, it treats all three as
\textbf{sources of variance} and measures how much each one contributes to the
spread in reported performance.


\section{Limitations of the Conventional Approach}

The standard approach to sepsis prediction follows a simple formula: choose one
Sepsis-3 definition, use one train/test split, report AUROC, and compare models.
Most of the ninety-one models reviewed by Wang et al.\ \cite{wang2025srpm}
follow this pattern.

This approach has four key problems:

\begin{enumerate}
    \item \textbf{The sepsis label is not fixed.} Cohen et al.\
    \cite{cohen2024subtle} showed that different ways of applying Sepsis-3 to
    the same database can create very different patient groups (867--2,178 cases).
    The label choice can change results as much as the model choice.

    \item \textbf{The data split can inflate results.} Guo et al.\
    \cite{guo2022temporal} found that sepsis prediction suffers the largest
    performance drop of any clinical task when tested on data from a different time
    period. Random splits let future patterns leak into training, making models
    look better than they are.

    \item \textbf{AUROC hides the real failure.} Wang et al.\
    \cite{wang2025srpm} showed that a model can keep a high AUROC while its
    clinical usefulness drops to harmful levels. The failure is invisible under
    AUROC alone.

    \item \textbf{The label captures doctor behaviour.} Kamran et al.\
    \cite{kamran2024evaluation} showed that models may learn when doctors
    suspect sepsis and begin treatment, not when the patient's condition actually
    worsens.
\end{enumerate}

A study that fixes all of these choices and compares only models is measuring
one source of variation while ignoring three others of similar size. This thesis
fills in the missing cells of the comparison.


\section{Research Question and Pre-Registered Hypotheses}
\label{sec:rq}

\noindent\textbf{Research question.} \textit{In real-time sepsis prediction on
MIMIC-IV, how much of the reported performance variation comes from (a) the
choice of model, and (b) the researcher's decisions: how sepsis is defined,
    how data is split, when prediction is made, and how performance is measured?}

Six hypotheses were pre-registered before any modelling:

\begin{description}
    \item[H1 (Label dominance)] Changing the sepsis definition affects AUROC at
    least as much as changing the model.
    \item[H2 (Coefficient instability)] At least one clinical feature changes its
    effect direction across label variants.
    \item[H3 (Treatment leakage)] Restricting evaluation to before treatment
    begins reduces AUROC by at least 0.03.
    \item[H4 (Split optimism)] Time-based testing gives at least 0.02 lower AUROC
    than random testing.
    \item[H5 (Metric divergence)] The Utility Score declines more than AUROC when
    moving from internal to external validation.
    \item[H6 (Model-class null)] The primary statistical model performs within 0.02
    AUROC of gradient-boosted trees.
\end{description}

H1 and H6 are in deliberate tension: if H1 is confirmed, the label is the main
source of variation; if H6 is rejected, the model also matters. Both outcomes
are informative.

These six form a single confirmatory family and are tested with
Holm--Bonferroni control of the family-wise error rate. Every other analysis in
this thesis, including the entire equity analysis, is exploratory and is
reported with cluster-bootstrap intervals and Benjamini--Hochberg
false-discovery control. Section~\ref{sec:inference} sets out the plan and
states plainly that it was written after the analyses had been run.


\section{Contributions}
\label{sec:contrib}

\begin{enumerate}
    \item \textbf{C1 (Label variance).} Label sensitivity itself is established
    by Cohen et al.\ \cite{cohen2024subtle} and Dutta et al.\
    \cite{dutta2026performance}, and neither the survival setting nor the
    multi-variant design is new. What is new, to the best of our knowledge, is
    the combination: three Sepsis-3 readings applied to a discrete-time
    competing-risks model on MIMIC-IV, with the label spread and the
    model-class spread measured on the same test rows and each carrying an
    interval, so that the two can be compared as estimands rather than as
    ranges of point estimates (Section~\ref{sec:search_strategy} states the
    search strategy and the bound on this claim).

    \item \textbf{C2 (Treatment leakage located in the onset rule).} A proof that
    a conjunction-based Sepsis-3 onset rule, the rule used here and widely in
    the MIMIC-IV literature, makes pre-treatment prediction structurally
    impossible, together with the measurement that the standard
    $\min(t_{\text{susp}}, t_{\text{SOFA}})$ rule applied to the identical stays
    recovers \pcPreTreatOnsetsVarF{} pre-treatment onsets at
    \pcPreTreatAurocPrimaryF{} AUROC. Kamran et al.'s
    \cite{kamran2024evaluation} concern is confirmed for label \emph{membership}
    and narrowed for label \emph{timing}.

    \item \textbf{C3 (Utility ceiling documented).} Swept across the whole
    threshold grid rather than scored at a conventional cut-off, no model
    exceeds \pcUtilityMaxGbtB{} under the primary label or
    \pcUtilityMaxGbtC{} under any pre-registered variant (a few per cent
    above the strategy of issuing no alerts, on a scale whose upper bound is 1), and reaching even that costs tens of false alerts per true one. This
    extends Wang et al.'s \cite{wang2025srpm} meta-analytic finding to a
    controlled single-study comparison and settles it against the
    operating-point objection.

    \item \textbf{C4 (Coefficient instability).} \pcHTwoEst{} of the
    \pcHTwoNTerms{} shared clinical covariates are unstable across label
    variants, \pcHTwoNSign{} of them by direction (including vasopressor
    exposure, a treatment indicator), supporting Lauritsen et al.'s
    \cite{lauritsen2021framing} framing-instability finding on MIMIC-IV.

    \item \textbf{C5 (Equity-by-label analysis).} Subgroup performance
    disparity and subgroup variation in who gets labelled are both reported by
    Wang et al.\ \cite{wang2022equity}, though for mortality prediction among
    patients already identified as septic rather than for onset prediction.
    What is new, to the best of our knowledge, is measuring the two on the
    \emph{same} strata in the same onset-prediction cohort under common
    false-discovery control, which is what shows that only one of them is
    precisely enough estimated to support a claim
    (Section~\ref{sec:search_strategy}).

    \item \textbf{C6 (Reproducible pipeline).} A pre-registered, seed-fixed
    pipeline following TRIPOD+AI \cite{collins2024tripodai} and PROBAST+AI
    \cite{moons2025probastai} reporting standards.

    \item \textbf{C7 (Definitional vs.\ empirical label variance separated).}
    The Sepsis-3 conjunction is split into its two limbs and each is labelled
    separately, so that the part of the label effect guaranteed by the
    definition can be told apart from the part discovered in the data. Two
    results follow. The onset \emph{time} under every pre-registered Sepsis-3
    variant is fixed entirely by the treatment anchor, the organ-dysfunction
    criterion filters which stays are labelled and moves no onset by a single
    hour, while a treatment-independent deterioration label agrees with
    Sepsis-3 at $\kappa = \pcKappaVsBD$, close to chance. And the treatment
    anchor alone is as predictable as the Sepsis-3 conjunction built on top of
    it ($\phi_{\text{anchor}} = \pcPhiAnchorPrimary$, interval including 1), so
    the difficulty lies neither in the conjunction nor in the physiology but in
    when the label is stamped. The same
    comparison shows the empty pre-treatment window of C2 to be a property of
    the onset rule, not of the patients: a treatment-independent label records
    \pcPreTreatOnsetsVarD{} onsets there at \pcPreTreatAurocPrimaryD{} AUROC,
    and Sepsis-3 itself records \pcPreTreatOnsetsVarF{} once its onset is timed
    as the standard specifies (Section~\ref{sec:label_anchor_results}).
\end{enumerate}


\section{Thesis Structure}

Chapter~\ref{ch:review} reviews the evidence that motivated this study's design.
Chapter~\ref{ch:methodology} describes the data, labels, models, and evaluation
plan.
Chapter~\ref{ch:design} explains the pipeline architecture and its twelve
scripts.
Chapter~\ref{ch:evaluation} presents all results, organised around the six
hypotheses.
Chapter~\ref{ch:discussion} interprets them, bounds what each estimate can
support, and states the limitations.
Chapter~\ref{ch:conclusion} draws conclusions and outlines future work.
Appendix~\ref{app:code} gives the source-code repository and the procedure for
reproducing every reported number.
